\documentclass[11pt,a4paper]{article}
% Estilo comum dos PDFs individuais da Sprint 1.
\usepackage[utf8]{inputenc}
\usepackage[T1]{fontenc}
\usepackage[brazilian]{babel}
\usepackage[a4paper,margin=2.35cm]{geometry}
\usepackage{xcolor}
\usepackage{hyperref}
\usepackage{enumitem}
\usepackage{listings}
\usepackage{parskip}
\usepackage{needspace}
\definecolor{azul}{RGB}{20,63,112}
\hypersetup{colorlinks=true,linkcolor=azul,urlcolor=azul}
\setlist[itemize]{leftmargin=*,itemsep=3pt}
\lstset{basicstyle=\ttfamily\small,frame=single,breaklines=true,columns=fullflexible}
\sloppy
\newcommand{\codigo}[1]{\texttt{#1}}
\newcommand{\titulo}[2]{\begin{center}{\LARGE\bfseries #1\par}\vspace{4pt}{\large #2\par}\vspace{8pt}\rule{0.8\textwidth}{0.4pt}\par\vspace{5pt}{\small Matheus Moreli Durão Silva -- Matrícula 12621CBS217\par Cibersegurança - Bacharelado - Noturno\par}\end{center}}
\newcommand{\cardfonte}[1]{\par\smallskip\noindent\textit{Referência: }\href{#1}{Trello - enunciado e checklist do card}.}

\begin{document}\titulo{S1\_08 - Estruturas de repetição}{Manter um menu aberto até ``Sair''}
\section*{Entregável correto}
O card pede um menu que se repita até a opção ``Sair''. Isso prepara a estrutura principal do sistema: o programa não realiza uma única ação e fecha; ele apresenta escolhas, executa uma ação e retorna ao menu.
\section*{Como cada laço funciona}
\codigo{while} repete enquanto uma condição for verdadeira. Um \codigo{while True} é útil para menu, desde que exista uma saída explícita com \codigo{break}. \codigo{for} percorre itens conhecidos; \codigo{range(1, 11)} produz números de 1 a 10. \codigo{continue} pula apenas a volta atual do laço, enquanto \codigo{break} encerra o laço inteiro.
\begin{lstlisting}
while True:
    opcao = input("1-Cadastrar | 0-Sair: ")
    if opcao == "0":
        print("Programa encerrado.")
        break
    elif opcao == "1":
        print("Cadastro selecionado.")
    else:
        print("Opcao invalida.")
        continue
\end{lstlisting}
\section*{Aplicação no projeto}
O menu do \codigo{app.py} aplica exatamente essa ideia: fica ativo, lê uma escolha, chama a função correspondente e termina somente na opção de saída. Laços \codigo{for} também listam ativos e vulnerabilidades um por um.
\section*{Checklist demonstrado}
Repetição, \codigo{for}, \codigo{while}, \codigo{range}, \codigo{break}, \codigo{continue} e aplicação prática foram estudados. Teste: selecionar uma opção, voltar ao menu e só então digitar ``0''; o encerramento não deve ocorrer antes.
\cardfonte{https://trello.com/c/op3s5sKj}
\end{document}
